\section{Marco teórico}    

\subsection{Norma IEC 60601}
    La norma IEC-60601 es una serie de estándares internacionales emitida por la Comisión Electrotécnica Internacional que
especifica los requisitos generales para la seguridad básica y el desempeño esencial de los equipos electromédicos (EM). 
Esta norma describe mecanismos rigurosos para proteger al paciente, al operador y al entorno de riesgos eléctricos (choques,
corrientes de fuga), mecánicos, térmicos, radiaciones no deseadas y fallas funcionales. Se relaciona directamente con
la norma ISO 14971 (Gestión de Riesgos en Dispositivos Médicos), la cual exige un análisis integral de riesgos durante todo el
ciclo de vida del producto para garantizar que los beneficios clínicos superen cualquier riesgo residual.

La estructura de la IEC-60601 se divide en una norma general (IEC 60601-1), normas colaterales (serie IEC 60601-1-X, sobre
compatibilidad electromagnética, alarmas, usabilidad) y normas específicas (serie IEC 60601-2-X para equipos particulares
como electrocardiógrafos o desfibriladores). En la práctica clínica y de ingeniería biomédica, se aplica en las etapas
de diseño, certificación, recepción de equipamiento, mantenimiento preventivo e inspección periódica en instituciones de salud
para validar que los equipos operen dentro de tolerancias seguras

\subsection{Seguridad y eficacia de equipos médicos}
Para que un equipo médico pueda utilizarse de manera confiable en un entorno clínico, 
debe cumplir dos condiciones fundamentales: ser seguro y ser eficaz.

La seguridad está relacionada con la capacidad del equipo para evitar daños al paciente y al operador durante 
su utilización. Esto incluye la protección frente a descargas eléctricas, corrientes de fuga, problemas de aislamiento
y otros riesgos que puedan surgir durante el funcionamiento normal o ante una condición de falla.
La eficacia, en cambio, se relaciona con la capacidad del equipo para cumplir correctamente con la función para la
 cual fue diseñado. En el caso de un monitor, por ejemplo, implica que las variables fisiológicas que mide se encuentren
 suficientemente próximas a los valores reales o de referencia y que sus resultados sean reproducibles. 
Ambos aspectos son importantes y están relacionados entre sí. Un equipo puede presentar buenas mediciones,
 pero si existe un riesgo eléctrico para el paciente, no puede considerarse adecuado para su utilización clínica. 
 Del mismo modo, un equipo eléctricamente seguro que entregue mediciones incorrectas tampoco cumpliría correctamente
   su función.
\subsection{Parte aplicable}

\subsection{Analizadores}